\chapter{Glossaire et révision}
\section*{API}
Contrat d’échange. Utilité et exemple : /api/resources relie interface et catalogue.

\section*{Middleware}
Étape de traitement HTTP. Utilité et exemple : requireUser refuse un rôle insuffisant.

\section*{DTO}
Représentation de données pour un échange. Utilité et exemple : serialize omet les clés privées.

\section*{ORM}
Correspondance objets/tables. Utilité et exemple : Prisma fournit findMany.

\section*{Transaction}
Modifications atomiques. Utilité et exemple : Compte et jetons sont créés ensemble.

\section*{Index}
Structure accélérant une recherche. Utilité et exemple : Le catalogue indexe ses critères.

\section*{Clé étrangère}
Référence garantie par la base. Utilité et exemple : subjectId référence Subject.

\section*{JWT}
Assertions signées. Utilité et exemple : La gateway vérifie rôle et signature.

\section*{Hash}
Empreinte irréversible. Utilité et exemple : SHA-256 protège le refresh token stocké.

\section*{Sel}
Aléa du hachage de mot de passe. Utilité et exemple : Bcrypt produit des hashes distincts pour un même secret.

\section*{CSRF}
Requête indésirable utilisant la session. Utilité et exemple : Origin est vérifié sur les mutations.

\section*{XSS}
Exécution de contenu injecté. Utilité et exemple : MathFormula garde un fallback texte.

\section*{CORS}
Règles navigateur interorigines. Utilité et exemple : La gateway déclare les origines.

\section*{Sandbox}
Environnement à capacités limitées. Utilité et exemple : Le convertisseur n’a pas de réseau.

\section*{Tmpfs}
Espace temporaire en mémoire. Utilité et exemple : La conversion dispose de 256 Mo.

\section*{Cache}
Réutilisation d’un résultat. Utilité et exemple : outputKey désigne le PDF conservé.

\section*{Idempotence}
Même effet après répétition. Utilité et exemple : Les seeds utilisent des upserts.

\section*{Révision}
Version d’un état métier. Utilité et exemple : revision refuse un résultat obsolète.

\section*{Hydratation}
Activation client du rendu serveur. Utilité et exemple : React connecte les boutons du dashboard.

\section*{État vide}
Vue d’absence de résultat. Utilité et exemple : Le catalogue invite à changer les filtres.

\section*{Streak}
Série de jours actifs. Utilité et exemple : Le résumé compare le dernier jour au jour courant.

\section*{Polling}
Interrogation périodique. Utilité et exemple : ResourceReader attend READY.

\section{Questions à résoudre dans le code}
Pourquoi vérifier les droits sur le fichier si sa carte était protégée ? Pourquoi le temps de travail ne mesure-t-il pas la maîtrise ? Pourquoi un document publié peut-il être en préparation ? Pourquoi une transaction SQL ne supprime-t-elle pas un fichier orphelin ? Retrouver la fonction puis proposer un test de refus ou de reprise pour chaque cas.
